% !TeX root = ../navier-stokes-manuscript.tex
% ============================================================
% 2.5 Simultaneous Sobolev Intersections
% ============================================================

\subsection{Simultaneous Sobolev Intersections}\label{sub:SimultaneousSobolevIntersections}

The spatial column becomes smooth one requested derivative at a time.
In three dimensions,
\[
  H^m(\R^3)
  \hookrightarrow
  C^k(\R^3)
  \qquad
  \text{whenever }
  m>k+\frac32.
\]
Fix one derivative order \(k\).
One finite Sobolev height \(m>k+3/2\) controls every spatial derivative through order \(k\).
Since the choice can be made for every finite \(k\),
\[
  \bigcap_{m=0}^{\infty}H^m(\R^3)
  \hookrightarrow
  C^\infty(\R^3).
\]
Each requested derivative uses one finite row high enough for that request.

The temporal column works in the same coordinatewise way on a strict finite interval \(I=(0,T)\), \(T<T_*\).
For a Hilbert space \(X\),
\[
  H^r(I;X)
  \hookrightarrow
  C^q(\overline I;X)
  \qquad
  \text{whenever }
  r>q+\frac12.
\]
Fixing \(q\) and choosing one finite \(r>q+1/2\) gives the desired temporal derivative.
Thus
\[
  \bigcap_{r=0}^{\infty}H^r(I;X)
  \hookrightarrow
  C^\infty(\overline I;X).
\]
Applied to the scalar critical height, this says
\[
  H\in
  \bigcap_{r=0}^{\infty}H^r(I)
  \quad\Longrightarrow\quad
  H\in C^\infty(\overline I).
\]
This scalar intersection describes the complete time regularity of \(H\).
The velocity field also carries spatial direction, phase, and geometry, all of which lie outside the scalar value \(H(t)\).

The velocity needs both indices.
Define its coordinatewise mixed intersection by
\[
  \mathscr V_\infty(I)
  :=
  \bigcap_{r=0}^{\infty}
  \bigcap_{m=0}^{\infty}
  H^r\bigl(I;H^m(\R^3)\bigr).
\]
The meaning of the intersection becomes concrete when one mixed derivative is requested.
For example, continuity of
\[
  \partial_t^2\partial_x^\alpha u,
  \qquad
  |\alpha|\leq4,
\]
follows from the single finite membership
\[
  u\in H^3\bigl(I;H^6(\R^3)\bigr),
\]
because \(3>2+1/2\) and \(6>4+3/2\).
In general, for any finite temporal order \(q\) and spatial order \(k\), choose finite indices
\[
  r>q+\frac12,
  \qquad
  m>k+\frac32.
\]
Then
\[
  H^r\bigl(I;H^m(\R^3)\bigr)
  \hookrightarrow
  C^q\bigl(\overline I;C^k(\R^3)\bigr).
\]
Arbitrary finite \(q\) and \(k\) give
\[
  \mathscr V_\infty(I)
  \hookrightarrow
  C^\infty\bigl(\overline I\times\R^3\bigr).
\]

Membership is asserted one finite coordinate at a time:
\[
  u\in\mathscr V_\infty(I)
  \quad\Longleftrightarrow\quad
  \norm{u}_{H^r(I;H^m)}
  <
  \infty
  \quad
  \text{for every finite }(r,m).
\]
The bound may change with \(r\) and \(m\).
The intersection places all of those finite statements on the same function \(u\).

The intersection \(\mathscr V_\infty(I)\) controls the velocity, but the Navier--Stokes history must also retain \(Q_n=\partial_t^np\), the differentiated momentum equation, the divergence constraint, and the pressure Poisson equation at every finite depth.
Even with those relations attached, a list of Sobolev memberships does not yet give a common time trace.
For a finite group of coordinates to reach the same time slice, adjacent spatial levels must meet around one middle space.
That finite group is the first rectangle.
